\subsection{What the retrieval interface changes}
\label{sec:actual-api}
The HeatpumpMonitor API permits full-range retrieval, so we tested it
directly \citep{hpmapi,hpmcode}. We fixed the first eligible
daily-persistence--MLP task from each external system before querying its
optional answers. Two repetitions of five policies produced 30 trials
and 60 HTTP requests, including common endpoint reads.

\begin{table}[htbp]\centering\small
\caption{Observed public-API control: three tasks, two repetitions per
policy. Incremental I/O excludes client throttle waiting. The full-range
response contains seven hourly samples, including five internal readings.
All five tested policies met the precision
target; the full five-policy measurements and cached response hashes are
in Online Resource~1, S15.}\label{tab:actual-api}
\begin{tabular}{lrrr}\toprule
Policy & Requests & Mean I/O (s) & Mean body bytes\\\midrule
Exact adaptive & 2 & 0.497 & 104.3\\
Full hourly range & 1 & 0.331 & 226.7\\\bottomrule
\end{tabular}\end{table}

Bulk retrieval uses one request, achieves zero terminal width and has
lower observed I/O time on these tasks. We therefore recommend bulk when
the full response is admissible. For selective access, the minimum-budget
calculation identifies the readings needed for assurance and the possible
reduction from feedback. Online Resource~1, S8 supplies the protocol selector
and its reading and acquisition-round cost catalogue.
